\documentclass[10pt,a4paper]{amsart}
\usepackage[margin=19mm]{geometry}
\usepackage{amsmath,amssymb,mathtools}
\usepackage[scr=rsfs,cal=euler]{mathalfa}
\usepackage{xfrac}
\usepackage[unicode,hidelinks,bookmarks=false]{hyperref}
\newcommand{\ie}{i.e.\ }
\newcommand{\cf}{cf.\ }
\newcommand{\resp}{respectively\ }
\newcommand{\Subsection}{\subsection}
\providecommand{\nobreakdash}{\nobreak\hbox{-}}
\setlength{\emergencystretch}{2em}
\multlinegap0pt
\usepackage{mathtools}
\usepackage{dsfont}
\usepackage{amssymb}
\newtheorem{thm}{Theorem}
\newcommand{\C}{\mathscr{C}}
\newcommand{\N}{\mathbb{N}}
\newcommand{\R}{\mathbb{R}}
\newcommand{\dd}{\mathrm{d}}
\newcommand{\eps}{\varepsilon}
\newcommand{\extendedR}{\R \cup \{+\infty\}}
\newcommand{\pr}{{\rm pr}}
\newcommand{\supp}{\textsf{supp}}
\title{A land of monotone plenty, \textit{bis repetita}: from classical to weak optimal transport}
\author{Virginie Ehrlacher, Rodrigue Lelotte, Luca Nenna}
\date{Source: arXiv:2606.19516v1 (CC BY 4.0)}
\begin{document}
\maketitle

\noindent\emph{Source and licence.} A land of monotone plenty, \textit{bis repetita}: from classical to weak optimal transport. Version arXiv:2606.19516v1 (CC BY 4.0), distributed by arXiv under the Creative Commons Attribution 4.0 International licence, http://creativecommons.org/licenses/by/4.0/, as recorded in the arXiv licence metadata for this exact version and shown on the abstract page https://arxiv.org/abs/2606.19516v1. Full licence notice: \texttt{licenses/arxiv-2606.19516-CC-BY-4.0.txt}.\par\smallskip

\noindent\textit{Editorial scope.} Counted unit: the article's thm thm:nec-cmon-bar (source0.tex lines 758--764). The article's complete original proof of it is delivered with it (source0.tex lines 766--881, one \texttt{proof} environment of 116 lines, which the article places away from the statement). No mathematics is written, repaired or re-derived: the statement and the proof are the article's own lines, in the article's own notation, and no retained line is substituted or re-flowed.  Retained uncounted as the source-local context the proof needs: lines 237--239; lines 421--423; lines 754--756.  Nothing was omitted from any retained span, so the package carries no omission record.  No retained line contains a citation, so the wrapper carries no bibliography and no bibliography entry.  No retained line contains a figure environment, an image-inclusion command or drawing code, so no image is delivered and the package declares no asset dependency.  Editorial formatting only, in addition: an AMS \texttt{amsart} preamble that restates the article's own declarations and macros used by the retained lines, namely the environments \texttt{thm} on separate counters, together with the standard packages those lines need.  The retained environments print in this document as Theorem 1 in order of appearance, whereas the article prints the same items with the numbering of its own full document.  The complete native source of the article as distributed in the arXiv e-print for this exact version is bundled unchanged as \texttt{sources/203152/source0.tex}.
\begin{thm}\label{thm:best-thm-necessary}
	Let $c : X \times Y \to \extendedR$ be Borel-measurable so that $c$ is (locally) integrable with respect to $\mu \otimes \nu$. If $\pi \in \Pi(\mu, \nu)$ is an optimal transport plan with $\C(\pi) < \infty$, then $\pi$ is $c$-cyclically monotone.
\end{thm}

\begin{equation}\label{WOT}\tag{WOT}
	\inf_{\pi \in \Pi(\mu, \nu)} \int_{X} C(x, \pi_x) \mu(\dd x)
\end{equation}

	\begin{multline}\label{eq:eq33_duh}
		-\int_{A \times \R^d} \left(\frac{x - \int_{\R^d} z \pi_x(\dd z)}{\Vert x - \int_{\R^d} z \pi_x(\dd z) \Vert}\right)^\intercal y [\gamma - \pi](\dd x, \dd y) \geq\\ - \int_{B_\gamma} \left\Vert \int_{\R^d} z \gamma_x(\dd z) - x\right\Vert \mu(\dd x)
	\end{multline}

	\begin{thm}\label{thm:nec-cmon-bar}
		Let $\pi \in \Pi(\mu, \nu)$ be an optimal transport plan for the weak optimal transport \eqref{WOT} with barycentric cost, \textit{i.e.} $C(x, \rho)= \Vert x - \int_{\R^d} z \rho(\dd z)\Vert$. Let us also assume that $\int_{\R^d} \Vert y \Vert \nu(\dd y) < \infty$. Then, there exists a measurable set $\Gamma \subset \R^d \times \R^d$ such that $\pi(\Gamma) = 1$ and so that, for all $n \in \N$ and all points $(x_1, y_1), \dots, (x_n, y_n) \in \Gamma$, we have 
		\begin{equation}\label{eq:c-cmon-bar}
		\boxed{\sum_{i | x_i \neq b_i} \frac{(y_{i+1} - y_i)^\intercal(b_i - x_i)}{\Vert b_i - x_i \Vert} \geq -\sum_{i | x_i = b_i} \Vert y_{i+1} -y_{i} \Vert}
		\end{equation}
	where we let $b_i := \int_{\R^d} z\pi_{x_i}(\dd z)$ and where we let $y_{n+1} := y_1$. 
	\end{thm}

	\begin{proof}[Proof of Theorem~\ref{thm:nec-cmon-bar}]
		Let $(x_1, y_1), \dots, (x_n, y_n) \in \Gamma$ where $\Gamma \subset X \times Y$ is a Borel set so that $\pi(\Gamma) = 1$. This set is to defined at the end of the proof by successive refinements, similar to the proof of Theorem~\ref{thm:best-thm-necessary}. For each $i = 1, \dots,n$, let $B_i^\eps \subset \R^d \times \R^d$ the ``ball'' of radius $\eps > 0$ centred at $(x_i, y_i) \in \R^d \times \R^d$. More precisely, to simplify a bit the presentation, assume that $B_i^\eps:=B_\eps(x_i) \times B_\eps(y_i)$ where $B_\eps(x) := \{ z \in \R^d : \Vert z - x \Vert < \eps\}$. We may (and will) assume without loss of generality --- up to refining $\Gamma$ ---  that $\pi(B_i^\eps) > 0$ for all $\eps > 0$ and all $i =1,\dots, n$. Let $\sigma_i^\eps$ be the restriction of $\pi$ on $B_i^\eps$ normalised to a probability measure, so that $\sigma_i^\eps(A) := \pi(B_i^\eps)^{-1}\pi(A \cap B_i^\eps)$ for all Borel set $A \subset \R^d \times \R^d$, and let $\sigma^\eps := \sum_{i = 1}^n \sigma_i^\eps$. We then define $\xi_i^\eps  := \pr_1^\sharp \sigma_i^\eps \otimes \pr_2^\sharp \sigma^\eps_{i+1}$ and $\xi^\eps := \sum_{i =1}^n \xi_i^\eps$. Letting $\gamma^\eps := \pi + t(\xi^\eps - \sigma^\eps)$, then $\gamma^\eps \in \Pi(\mu, \nu)$ for all $\eps > 0$ and any $0 < t < t_\eps$ where $t_\eps := \Vert \frac{d \sigma^\eps}{d \pi} \Vert_{L^\infty(\pi)}$. We apply the inequality \eqref{eq:eq33_duh} to $\gamma \leftarrow \gamma^\eps$, which gives
	 \begin{multline}\label{eq:eq35}
	 	-\int_{A \times \R^d} \frac{(x - \int_{\R^d} z \pi_x(\dd z))^\top y}{\Vert x - \int_{\R^d} z \pi_x(\dd z) \Vert} [\gamma^\eps - \pi](\dd x, \dd y) \geq \\- \int_{B_{\gamma^\eps}} \left\Vert \int_{\R^d} z \gamma_x^\eps(\dd z) - x\right\Vert \mu(\dd x)
	 \end{multline}
	 where we recall that the measurable sets $A$ and $B_{\gamma^\eps}$ are defined by
     $$
     A = \left\{x \in \R^d : x \neq \int_{\R^d} z \pi_x(\dd z) \right\}
     $$
     and 
     \begin{equation}\label{eq:B}
     B_{\gamma^\eps} = \left\{ x \in A^c : x \neq \int_{\R^d} z \gamma^\eps(\dd z)\right\}.
     \end{equation}
     The above inequality \eqref{eq:eq35} further develops to
	 \begin{multline}\label{eq:eq36}
	 		-\int_{A \times \R^d} \frac{(x - \int_{\R^d} z \pi_x(\dd z))^\top y}{\Vert x - \int_{\R^d} z \pi_x(\dd z) \Vert} [\xi^\eps - \pi^\eps](\dd x, \dd y) \geq \\- \int_{B_{\gamma^\eps}} \left\Vert \int_{\R^d} z [\xi^\eps_x - \pi^\eps_x](\dd z) \right\Vert \mu(\dd x)
	 \end{multline} 
     Here, $(\xi^\eps_x)_{x \in \R^d}$ and $(\sigma^\eps_x)_{x \in \R^d}$ are the disintegrations of $\xi^\eps$ and $\sigma^\eps$ with respect to $\mu$. We have $\sigma^\eps_{x} = \sum_{i=1}^n \sigma^\eps_{i, x}$ where for all $i = 1, \dots, n$ the disintegration $(\sigma^\eps_{i, x})_{x \in X}$ of $\sigma^\eps_i$ with respect to $\mu$ reads for $x \in B_\eps(x_i)$
	 \begin{equation}
	 	\sigma^\eps_{i, x}(\dd y)= \frac{1}{\pi(B_i^\eps)}\pi_{x}\restriction_{B_\eps(y_i)}(\dd y)
	\end{equation}
	--- where $\pi_{x}\restriction_{B_\eps(y_i)}$ denotes the restriction of $\pi_x$ to the ball $B_\eps(y_i)$ --- and $\sigma^{\eps}_{i, x} = 0$ for $x \notin B_\eps(x_i)$ otherwise. Likewise, $\xi^\eps = \sum_{i=1}^n \xi^\eps_{i, x}$ where the disintegration $(\xi^\eps_{i, x})_{x \in X}$ of $\xi^\eps_i$ with respect to $\mu$ reads for $x \in B_\eps(x_i)$
	\begin{equation}
		\xi^\eps_{i, x}(\dd y) = \frac{1}{\pi(B_i^\eps)} \frac{\pi_x(B_\eps(y_i))}{\pi(B_{i+1}^\eps)} \int_{B_\eps(x_{i+1})} \pi\restriction_{\R^d \times B_\eps(y_{i+1})}(\dd x', y)
	\end{equation}
    and $\xi^\eps_{i, x} = 0$ for $x \notin B_\eps(x_i)$. To show that \eqref{eq:c-cmon-bar} holds, we are going to take the limit $\eps \to 0$ in the inequality \eqref{eq:eq36}.
    
	 First, we note that, by the standard \textit{Lebesgue differentiation theorem} and just like in the proof of Theorem~\ref{thm:best-thm-necessary}, it holds that --- up to refining the set $\Gamma$:
	 \begin{multline}\label{eq:ldt-simple}
    -\int_{A \times \R^d} \frac{(x - \int_{\R^d} z \pi_x(\dd z))^\top y}{\Vert x - \int_{\R^d} z \pi_x(\dd z) \Vert} [\xi^\eps - \pi^\eps](\dd x, \dd y) \\ \xrightarrow[\eps \to 0]{} \sum_{i | x_i \neq b_i} \frac{(y_{i+1} - y_i)^\top (b_i - x_i)}{\Vert b_i - x_i \Vert}
	 \end{multline}
     where we recall the notation $b_i := \int_{\R^d} z \pi_{x_i}(\dd z)$ for all $i = 1, \dots, n$. This is exactly the left-hand side of the sought-for inequality \eqref{eq:c-cmon-bar}. Now, it remains to prove that we can recover the right-hand side of \eqref{eq:c-cmon-bar} from the right-hand side of \eqref{eq:eq36} in the limit $\eps \to 0$.
     
     Let us first remark that the right-hand side rewrites as 
	 \begin{multline}
	     \int_{B_{\gamma^\eps}} \left\Vert \int_{\R^d} z [\xi^\eps_x - \pi^\eps_x](\dd z) \right\Vert \mu(\dd x) \\ = \sum_{i = 1}^n 	\int_{B_{\gamma^\eps} \cap B_\eps(x_i)} \left\Vert \int_{\R^d} z [\xi^\eps_x - \pi^\eps_x](\dd z) \right\Vert \mu(\dd x)
	 \end{multline}
    Indeed, here we used that $x \mapsto \xi_x^\eps$ and $x \mapsto \pi_x^\eps$ are supported onto the balls $B_\eps(x_i)$'s. Using the explicit formula for the disintegrations given above, we have 
     \begin{multline}
             \int_{B_{\gamma^\eps} \cap B_\eps(x_i)} \left\Vert \int_{\R^d} z [\xi^\eps_x - \pi^\eps_x](\dd z) \right\Vert \mu(\dd x) =\\
             \frac{1}{\pi(B_i^\eps)} \int_{B_i^\eps} \mathds{1}_{B_{\gamma^\eps}}(x) \Bigg\Vert \frac{1}{\pi_x(B_\eps(y_i))}\int_{B_\eps(y_i)}z \pi_x(\dd z) \\- \frac1{\pi(B_{i+1}^\eps)} \int_{B_{i+1}^\eps} y\pi(\dd x', \dd y)\Bigg \Vert \pi(\dd x, \dd y)
     \end{multline}
     For all $\eps > 0$ (small enough) we have 
     \begin{multline}
     \Vert y_{i+1} - y_i \Vert - 2\eps \\\leq \left\Vert \frac{1}{\pi_x(B_\eps(y_i))}\int_{B_\eps(y_i)}z \pi_x(\dd z) - \frac1{\pi(B_{i+1}^\eps)} \int_{B_{i+1}^\eps} y\pi(\dd x', \dd y)\right \Vert \\\leq \Vert y_{i+1} - y_i \Vert + 2\eps
     \end{multline}
     and therefore 
     \begin{multline}
         (\Vert y_{i+1} - y_i \Vert - 2\eps)\frac{1}{\pi(B_i^\eps)} \int_{B_i^\eps} \mathds{1}_{B_{\gamma^\eps}}(x) \pi(\dd x, \dd y)  \\
         \leq
         \int_{B_{\gamma^\eps} \cap B_\eps(x_i)} \left\Vert \int_{\R^d} z [\xi^\eps_x - \pi^\eps_x](\dd z) \right\Vert \mu(\dd x) \\
         \leq (\Vert y_{i+1} - y_i \Vert + 2\eps)\frac{1}{\pi(B_i^\eps)} \int_{B_i^\eps} \mathds{1}_{B_{\gamma^\eps}}(x) \pi(\dd x, \dd y)
     \end{multline}


We now claim that there exists $\overline{\eps} > 0$ small enough so that we have the set inclusion $B_{\gamma^\eps} \subset B_{\gamma^{\overline{\eps}}}$ for all $\eps < \overline{\eps}$. Indeed, let us remark that the set $B_{\gamma^\eps}$ rewrites as the set of points $x \in X$ so that $x = \int_{\R^d} z \pi_x(\dd z)$ and so that $\int_{\R^d} z \xi^\eps_x(\dd z) \neq \int_{\R^d} z \sigma^\eps_x(\dd z)$. This last condition rewrites, using the disintegrations, as 
\begin{equation}\label{eq:b_gamma_eps_}
    \frac{1}{\pi_x(B_\eps(y_i))} \int_{B_\eps(y_i)} z \pi_x(\dd z) \neq \frac{1}{\pi(B_{i+1}^\eps)} \int_{B_{i+1}^\eps} y \pi(\dd x', \dd y)
\end{equation}

     We note that, as $\eps \to 0$, this condition reads (formally for the moment) $y_i \neq y_{i+1}$. Of course, we may assume without loss of generality that $y_i \neq y_{i+1}$. Indeed, if it were that $y_i = y_{i+1}$, then
\begin{equation}
    \int_{B_{\gamma^\eps} \cap B_\eps(x_i)} \left\Vert \int_{\R^d} z [\xi^\eps_x - \pi^\eps_x](\dd z) \right\Vert \mu(\dd x) \leq \Vert y_{i+1} - y_i \Vert + 2\eps = 2\eps
\end{equation}
and thus, in the limit $\eps \to 0$, this integral does not contribute to the sum since it vanishes to zero. Therefore, only the case where $y_{i+1} \neq y_i$ is of importance. We note passing by that the previous result implies trivially that, in the limit $\eps \to 0$, the inequality \eqref{eq:eq36} implies
\begin{equation}
\sum_{i | x_i \neq b_i} \frac{(y_{i+1} - y_i)^\intercal(b_i - x_i)}{\Vert b_i - x_i \Vert} \geq -\sum_{i = 1}^n \Vert y_{i+1} -y_{i} \Vert.
\end{equation}
But the sought-for inequality \eqref{eq:c-cmon-bar} is finer since only the indices $i \in \{1, \dots, n\}$ where $x_i = b_i$ are consider in the sum on the right-hand side. Let us therefore assume that $y_i \neq y_{i+1}$ for all $i = 1, \dots,n$. If $x \in B_{\gamma^{\eps}}$, and using that $y_{i+1} \neq y_i$, then there exists a small enough $\overline{\eps} > 0$ such that \eqref{eq:b_gamma_eps_} holds for all $0 < \eps < \overline{\eps}$. Indeed, this follows from the obvious facts that 
$$
\left\Vert\frac{1}{\pi_x(B_\eps(y_i))} \int_{B_\eps(y_i)} z \pi_x(\dd z) - y_i \right\Vert \leq \eps
$$
and 
$$
\left\Vert \frac{1}{\pi(B_{i+1}^\eps)} \int_{B_{i+1}^\eps} y \pi(\dd x', \dd y) -y_{i+1} \right \Vert \leq \eps
$$
for all $\eps > 0$. Therefore, letting for instance $\overline{\eps} < \frac{\Vert y_{i+1} - y_i \Vert}{2}$, we have the set inclusion $B_{\gamma^{\eps}} \subset B_{\gamma^{\overline{\eps}}}$ for all $0 < \eps < \overline{\eps}$. 

Using this fact, we obtain 
     \begin{multline}
         \limsup_{\eps \to 0}\int_{B_{\gamma^\eps} \cap B_\eps(x_i)} \left\Vert \int_{\R^d} z [\xi^\eps_x - \pi^\eps_x](\dd z) \right\Vert \mu(\dd x) \\\leq \limsup_{\eps \to 0} \frac{\Vert y_{i+1} - y_i\Vert + 2\eps}{\pi(B_i^\eps)} \int_{B_i^\eps} \mathds{1}_{B_{\gamma^{\overline{\eps}}}}(x) \pi(\dd x, \dd y)%\\=\Vert y_{i+1} - y_i\Vert \mathds{1}_{B_{\gamma^{\overline{\eps}}}}(x_i)
     \end{multline}
Now, according to the Lebesgue differentiation theorem, we know that  there exists a set $E \subset \R^d$ so that $\mu(E) = 1$ and so that 
     $$
     \frac1{B_\eps(x)} \int_{B_\eps(x)} \mathds{1}_{B_{\gamma^{\overline{\eps}}}}(x') \mu(\dd x') \xrightarrow[\eps \to 0]{} \mathds{1}_{B_{\gamma^{\overline{\eps}}}}(x)
     $$
     for all $x \in E$. \textbf{If we could ensure that $x_i \in E$ for all $i = 1, \dots,n$}, then we would obtain 
     $$
      \limsup_{\eps \to 0}\int_{B_{\gamma^\eps} \cap B_\eps(x_i)} \left\Vert \int_{\R^d} z [\xi^\eps_x - \pi^\eps_x](\dd z) \right\Vert \mu(\dd x) \leq \Vert y_{i+1} - y_i\Vert \mathds{1}_{B_{\gamma^{\overline{\eps}}}}(x_i) 
     $$
     and thus the claimed inequality \eqref{eq:c-cmon-bar} would immediately follows. Let us discuss how to ensure that this happens. 
     
     First, beware that the measurable set $E \equiv E[n, \mathcal{X}, \overline{\eps}]$ actually depends on the function $\mathds{1}_{B_{\gamma^{\overline{\eps}}}}$ which itself depends on $n$, the family of points $\mathcal{X} := \{(x_i, y_i) : i = 1, \dots, n\}$ under consideration, and $\overline{\eps} > 0$. To simplify the presentation, let us write $B_{\gamma^{\overline{\eps}}} = B[n, \mathcal{X}, \overline{\eps}]$ as well. The idea is then: since any subspace of a metrisable separable space is itself separable, the support $\supp(\pi)$ of the transport plan $\pi$  is separable. Let us then consider $S \subset \supp(\pi)$ to be a dense and countable subset of $\supp(\pi)$. For each $n \in \N$, the set $\mathbb{X}_n = \{\mathcal{X}\}$ of all families $\mathcal{X} = \{(x_1, y_1), \dots, (x_n, y_n) \} \subset S$ of $n$ elements inside $S$ is therefore itself a countable set as a finite product of countable sets. Then, for each $\overline{\eps} > 0$, consider the set 
$$
E^{\overline{\eps}} := \bigcap_{n \in \N} \bigcap_{\mathcal{X} \in \mathbb{X}_n} E[n, \mathcal{X}, \overline{\eps}]
$$
The set $E^{\overline{\eps}}$ is measurable and $\mu(E^{\overline{\eps}}) = 1$. Then, we finally consider 
$$
E := \bigcap_{m \in \N} E^{1/m}
$$ 
The set $E$ is again measurable and $\mu(E) = 1$. We then consider $\Gamma := (E \times \R^d)\cap \Gamma_0$ where $\Gamma_0$ is the $\pi$-full measure set over which the convergence ensured by the Lebesgue differentiation theorem holds in \eqref{eq:ldt-simple}. We claim that the sought-for inequality \eqref{eq:c-cmon-bar} holds over $\Gamma$. This follows from the fact, for any $n \in \N$ and any family 
$$
\mathcal{X} = \{(x_1, y_1), \dots, (x_n, y_n) \} \subset \Gamma \subset (E \times \R^d),$$
(where \textit{a priori} $\mathcal{X} \not\subset S$) and for all $\eta > 0$, there exists a family $\mathcal{X}' \in \mathbb{X}_n$ which is $\eta$-close to $\mathcal{X}$, in the sense that, letting
$$
\mathcal{X}' = \{(x_1', y_1'), \dots, (x_n', y_n') \},
$$ 
we have $\Vert (x_i, y_i) - (x_i', y_i')\Vert < \eta$ for all $i = 1, \dots, n$. Now, let us consider $\eta = \eps/2$. Then, by similar arguments to that of preceding, we have that $B[n, \mathcal{X}, \eps] \subset B[n, \mathcal{X}', \eps/2]$. In particular, selecting $\overline{\eps} = 1/m$ for $m \in \N$ big enough, we obtain
 \begin{equation}
\frac1{\pi(B_i^\eps)} \int_{B_i^\eps} \mathds{1}_{B[n, \mathcal{X}, \overline{\eps}]}(x) \pi(\dd x, \dd y) \leq \frac1{\pi(B_i^\eps)} \int_{B_i^\eps} \mathds{1}_{B[n, \mathcal{X}', \overline{\eps}/2]}(x) \pi(\dd x, \dd y).
\end{equation}
Finally, the right-hand side converges to $\mathds{1}_{B[n, \mathcal{X}', \overline{\eps}/2]}(x_i) \Vert y_{i+1} - y_i\Vert$ since $x_i \in E$. This terminates the proof.

     
   
	\end{proof}

\end{document}
